\begin{corollary}[5.7.10]
\label{IV.2.s5.7.10-fr}
Supongamos verificadas las hipótesis de \textup{(5.7.9)}, y supongamos además que $A$
es un anillo local; entonces, para que $M$ sea un $A$-módulo de
Cohen-Macaulay, es necesario y suficiente que, para todos los ideales primos
$\mathfrak q_i$ de $B$ por encima del ideal maximal de $A$,
$M_{\mathfrak q_i}$ sea un $B_{\mathfrak q_i}$-módulo de Cohen-Macaulay y,
además, que todos los números
$\dim_{B_{\mathfrak q_i}}(M_{\mathfrak q_i})$ sean iguales.
\end{corollary}

En efecto, resulta de (5.7.9.1) y de (5.7.9.2) que estas condiciones
equivalen a la relación
$\dim_A(M)=\operatorname{prof}_A(M)$.

\begin{corollary}[5.7.11]
\label{IV.2.s5.7.11-fr}
Supongamos verificadas las hipótesis de \textup{(5.7.9)}.
\begin{enumerate}
\item[\textup{(i)}] Si $M_{[\rho]}$ verifica la propiedad $(S_k)$, lo mismo ocurre
con $M$.
\item[\textup{(ii)}] Supongamos que, para todo par de ideales primos
$\mathfrak q$, $\mathfrak q'$ de $B$ tales que
$\rho^{-1}(\mathfrak q)=\rho^{-1}(\mathfrak q')$, se tiene
$\dim_{B_{\mathfrak q}}(M_{\mathfrak q})
=\dim_{B_{\mathfrak q'}}(M_{\mathfrak q'})$. Entonces, si $M$ verifica la
propiedad $(S_k)$, lo mismo ocurre con $M_{[\rho]}$.
\end{enumerate}
\end{corollary}

\oldpage[IV-2]{107}

Esto resulta inmediatamente de las relaciones (5.7.9.1) y (5.7.9.2) y de la
definición de la propiedad $(S_k)$.

\begin{env}[5.7.12]
\label{IV.2.s5.7.12-fr}
De conformidad con las definiciones de (5.7.1), dados un anillo noetheriano
\emph{cualquiera} $A$ y un $A$-módulo $M$ de tipo finito, se define
$\operatorname{coprof}_A(M)$ como igual a
$\operatorname{coprof}(M)=\sup_{x\in X}
(\operatorname{coprof}_{A_x}(M_x))$, donde $X=\operatorname{Spec}(A)$;
veremos más adelante (6.11.5) que esta definición \emph{coincide} con la de
(\textbf 0, 16.4.9) cuando $A$ es un anillo noetheriano \emph{local}.
\end{env}

\begin{corollary}[5.7.13]
\label{IV.2.s5.7.13-fr}
Sean $A$, $B$ dos anillos noetherianos, $\rho:A\to B$ un homomorfismo
de anillos y $M$ un $B$-módulo tal que $M_{[\rho]}$ sea un $A$-módulo de tipo
finito. Entonces se tiene
\begin{equation}
\operatorname{coprof}_A(M_{[\rho]})\geq\operatorname{coprof}_B(M).
\tag{5.7.13.1}\label{IV.2.s5.7.13.1-fr}
\end{equation}
\end{corollary}

Esto resulta de la definición precedente y de las relaciones (5.7.9.1) y
(5.7.9.2).
